\section{Discussion, Cost, and Limitations}
\subsection{Causal and Temporal Interpretation}
The detector supplies a warning, PGD-2 restores robustness, and
immediate-switch converts the warning into
timely action so that PGD-2 can begin before CO. Immediate-switch matches
immediate-replay's 5/5 pre-collapse interventions and differs by only $+0.22$ PGD-10
points, with a paired interval crossing zero. Restoring the epoch snapshot is
therefore not required for the observed timing benefit.

Table~\ref{tab:event-audit} reports every matched timing event. High lead time
is desirable for safety, but a warning far ahead of collapse does not imply
precise prediction of collapse timing; the monitor acts as a safety trigger
upon observing geometry degradation. In particular, the 1568- and 1680-update
leads in seeds 23 and 202 should not be interpreted as calibrated time-to-CO
forecasts. The saved-trace audit in the Appendix also checks whether the score
remains above threshold after these early warnings.

\begin{table}[t]
\centering
\caption{Event-level timing in optimizer updates}
\label{tab:event-audit}
\scriptsize
\setlength{\tabcolsep}{3pt}
\begin{tabular}{rrrrc}
\toprule
Seed & Warn$\to$CO & Immediate-switch & Deferred & Deferred before CO \\
\midrule
17  & 40   & 0 & 0   & yes \\
23  & 1568 & 0 & 212 & yes \\
42  & 20   & 0 & 232 & no \\
101 & 60   & 0 & 292 & no \\
202 & 1680 & 0 & 52  & yes \\
\bottomrule
\end{tabular}
\end{table}

\subsection{Engineering Cost and Reproducibility}
The attack-start gradient constructs the FGSM example. The endpoint input
gradient is obtained from the same parameter-backward pass, without an
additional forward or backward pass. The monitor retains it only until its
cosine is reduced to a batch scalar, then stores scalar window history. It does
not call a validation attack or an extra model evaluation.

The primary immediate-switch response does not require a snapshot. Its 18.74-minute
measurement is still core training time, not complete end-to-end runtime. The
immediate-replay control has a separate systems cost: its implementation holds one
epoch-start copy of model, optimizer, and AMP-scaler state in device memory,
plus Python, NumPy, PyTorch CPU/CUDA, and loader-generator RNG state. It does
not write this immediate-replay snapshot to disk. Reported core training time does not
include snapshot save-restore overhead. It does include abandoned FGSM
parameter-update time and the replayed PGD-2 epoch. The experiment did not
instrument snapshot-copy latency or peak memory, so the current measurements
must not be read as complete end-to-end latency or memory results.

Every detector choice is fixed before its associated held-out evaluation:
20-update windows, five-window median history, three-window minimum history,
the event definition, and the 60-update near horizon. Fixed-switch controls
ignore detector callbacks. Retrospective PGD labels never trigger the proposed
monitor. Seed-level warning, collapse, PGD-start, accuracy, and timing records
remain in the saved experiment CSV files; paired intervals use 100,000 resamples.

\subsection{Scope and Limitations}
The architecture result is partial unchanged-threshold transfer: ResNet-18
detects 4/5 events but only 2/5 near events. CIFAR-100 is not zero-shot or
plug-and-play; transfer requires six dataset-specific development trajectories
to refit the detector operating point using its original event-level
lexicographic rule. Only one of its three collapse-prone development runs meets
the CO definition, so the event-level calibration has limited statistical
stability. As a supplementary check, a window-level audit contains one positive
next-window label among 2,877 eligible windows. One alarm episode
occurs across three held-out non-event runs; its unnecessary PGD-2 cost is
unquantified.
Event samples for ResNet-18, CIFAR-100, and the
$16/255$ stress test remain small, as reflected by wide Wilson intervals.
Transfer evidence is still concentrated in CIFAR datasets and the ResNet
family; ImageNet and substantially different architectures remain outside the
current scope.

The primary immediate-switch checkpoint is included in full-test PGD-50 and
multi-restart PGD, plus the full standard AutoAttack suite on a fixed
1,000-image subset, but not full-test
AutoAttack. The subset is identical across compared checkpoints and includes
APGD-CE, APGD-T, FAB-T, and Square. Agreement among PGD and this supplementary
evaluation reduces, but does not eliminate, concern about gradient masking.
The cosine score has empirical and mechanistic support but is neither a
theoretical certificate nor a necessary or sufficient condition for CO.
The $16/255$ recovery results use immediate-replay; immediate-switch has not
been tested in that stress setting. Immediate-replay also has snapshot memory,
storage-policy, and engineering overhead not characterized for larger models.
Immediate-switch avoids this snapshot and has both five-seed intervention and
pre-fixed seed-101 strong evaluations in Tables~\ref{tab:test}--\ref{tab:aa}.

GradAlign has higher mean PGD-10 robustness than the proposed policy, which uses lower
measured core training time and provides an explicit event log, but the choice
between them depends on whether the deployment prioritizes peak robustness,
measured training cost, or auditable failure detection.
